%% IEEE-formatted technical report (derived from Data 202 Report.txt)
%% Compile with: pdflatex Data_202_Report.tex (requires IEEEtran.cls)

\documentclass[conference]{IEEEtran}

\usepackage{cite}
\usepackage{amsmath,amssymb}
\usepackage{graphicx}
\usepackage{textcomp}
\usepackage{xcolor}
\usepackage[hidelinks]{hyperref}

\hyphenation{op-tical net-works semi-conduc-tor}

\begin{document}

\title{Statistical Analysis of User--Movie Interaction Data for Recommendation Systems}

\author{
\IEEEauthorblockN{Luca Popescu,
River Roseveare-Hunt,
Timothy Tran,
and Cliff Gonsalves}
\IEEEauthorblockA{\textit{DATA 202 Course Project}}
}

\maketitle

\begin{abstract}
This report analyzes user--movie interaction data to better understand patterns in user preferences and item popularity within a recommendation system. The dataset consists of over 3.4 million ratings from nearly 300{,}000 users across more than 60{,}000 items. We apply descriptive statistics, correlation analysis, confidence intervals, and hypothesis testing to evaluate rating behavior and structural characteristics of the dataset. Results show that ratings are heavily skewed toward higher values, with a mean rating of 4.22 and a median of 5. Additionally, the user--item matrix is extremely sparse, and user activity is highly uneven. Hypothesis testing confirms that the average rating is statistically greater than a benchmark value of 4.0, although the practical difference is small. Correlation analyses reveal weak relationships between rating behavior and time or user activity. These findings highlight key challenges in recommendation modeling, particularly sparsity and skewness, and inform future approaches such as collaborative filtering.
\end{abstract}

\begin{IEEEkeywords}
Recommendation systems, descriptive statistics, hypothesis testing, correlation analysis, data sparsity, collaborative filtering.
\end{IEEEkeywords}

\section{Introduction}

The goal of this project is to analyze user--movie interaction data within a recommendation system to understand the factors that influence user preferences and item popularity. Recommendation systems are widely used in platforms such as streaming services and e-commerce, making it important to understand the structure and behavior of rating data.

This analysis focuses on identifying patterns in user ratings, evaluating statistical properties of the dataset, and assessing relationships between variables such as time and user activity. Understanding these characteristics is essential for building effective recommendation models.

\section{Dataset Description}

The dataset consists of user ratings for movies and TV items collected over time.

\subsection{Key Dataset Characteristics}

\begin{itemize}
    \item Total observations: 3{,}410{,}019
    \item Unique users: 297{,}529
    \item Unique items: 60{,}175
    \item Rating scale: 1 to 5
    \item Time range: December 2, 1997 to October 1, 2018
\end{itemize}

\subsection{Variables}

\begin{itemize}
    \item \texttt{user\_id}: categorical identifier for users;
    \item \texttt{item\_id}: categorical identifier for items;
    \item \texttt{rating}: discrete numerical variable (1--5);
    \item \texttt{timestamp}: numeric time variable (converted to datetime).
\end{itemize}

\subsection{Data Preparation}

\begin{itemize}
    \item Verified dataset structure and column consistency;
    \item Confirmed correct data types for ratings and timestamps;
    \item Converted Unix timestamps into readable calendar dates;
    \item Computed summary statistics (e.g., ratings per user/item);
    \item Estimated sparsity of the user--item matrix.
\end{itemize}

The dataset was found to be highly sparse, meaning most user--item interactions are unobserved. This has important implications for modeling.

\section{Mathematical and Statistical Methods}

We applied several statistical methods to analyze the dataset.

\subsection{Descriptive Statistics}

Descriptive statistics summarize the rating distribution:
\begin{itemize}
    \item mean, median, mode;
    \item standard deviation;
    \item quartiles.
\end{itemize}

\subsection{Hypothesis Testing}

We conducted a one-sample $t$-test to evaluate whether the average rating differs from a benchmark value of 4.0:
\begin{align}
    H_0 &: \mu = 4.0, \\
    H_1 &: \mu \neq 4.0.
\end{align}

The test statistic is
\begin{equation}
    t = \frac{\bar{x} - \mu_0}{s / \sqrt{n}},
\end{equation}
where $\bar{x}$ is the sample mean, $\mu_0$ is the hypothesized mean, $s$ is the sample standard deviation, and $n$ is the sample size.

\subsection{Confidence Interval}

A 95\% confidence interval for the mean rating (large-sample normal approximation) is
\begin{equation}
    \bar{x} \pm z_{\alpha/2} \cdot \frac{s}{\sqrt{n}},
\end{equation}
where $z_{\alpha/2}$ is the critical value for a two-sided interval at level $\alpha = 0.05$. For finite samples, the $t$ distribution may be preferred.

\subsection{Correlation Analysis}

The Pearson correlation coefficient measures linear association:
\begin{equation}
    r = \frac{\operatorname{cov}(X,Y)}{\sigma_X \sigma_Y}.
\end{equation}

These methods provide both descriptive insights and inferential conclusions about the dataset.

\section{Analysis and Results}

\subsection{Rating Distribution}

\begin{itemize}
    \item Mean: 4.2213
    \item Median: 5
    \item Mode: 5
    \item Standard deviation: 1.17
    \item Quartiles: $Q_1 = 4$, $Q_2 = 5$, $Q_3 = 5$
\end{itemize}

\textbf{Findings:} Ratings are heavily concentrated at the high end; five-star ratings dominate; at least 75\% of ratings are 4 or higher.

\subsection{User and Item Structure}

\begin{itemize}
    \item Average ratings per user: 11.46 (median: 7)
    \item Average ratings per item: 56.67 (median: 16)
    \item Sparsity: 0.999810
\end{itemize}

\textbf{Findings:} User activity is highly uneven; item popularity follows a long-tail distribution; the dataset is extremely sparse.

\subsection{Hypothesis Testing Results}

\begin{itemize}
    \item Sample mean: 4.2213
    \item 95\% CI: $[4.2201,\, 4.2226]$
    \item $t$-statistic: 350.37
    \item $p$-value: $< 0.001$
\end{itemize}

\textbf{Interpretation:} Reject $H_0$; the mean rating is statistically greater than 4.0. The large sample size leads to extremely strong statistical significance.

\subsection{Correlation Analysis}

\textbf{Ratings vs.\ time:}
\begin{itemize}
    \item $r = 0.0938$
    \item $p$-value $= 0.6780$
\end{itemize}
$\rightarrow$ Very weak; not statistically meaningful at conventional levels for the stated test.

\textbf{User activity vs.\ rating behavior:}
\begin{itemize}
    \item $r = -0.0235$
    \item $p$-value $\approx 0$
\end{itemize}
$\rightarrow$ Statistically significant but practically negligible.

\section{Discussion}

The dataset reveals several important structural and behavioral patterns:
\begin{itemize}
    \item Ratings are strongly skewed toward higher values, indicating possible bias or user satisfaction trends.
    \item User engagement is uneven, with many users contributing few ratings.
    \item Item popularity follows a long-tail distribution: a small number of items receive most ratings.
    \item Extreme sparsity presents challenges for modeling.
\end{itemize}

Although hypothesis testing shows statistical significance, the practical difference in average rating is small. Similarly, correlations---while sometimes statistically significant---are too weak to be actionable.

These findings suggest that simple models may struggle, and more advanced techniques (e.g., collaborative filtering) are needed.

\section{Conclusion}

This analysis explored user--movie interaction data using statistical and mathematical methods. The dataset exhibits strong positive rating bias, uneven user behavior, and extreme sparsity. While statistical tests confirm that the average rating differs from a benchmark, the practical significance is limited.

These insights highlight key challenges for recommendation systems and motivate the use of more advanced modeling approaches in future work.

\section*{Acknowledgment}

The authors thank the DATA 202 course staff for course materials and guidance.

\begin{thebibliography}{00}

\bibitem{course} Course materials, DATA 202, Department course notes and slides.

\bibitem{amazon} J.~McAuley and collaborators, ``Amazon review data,'' UC San Diego. [Dataset for Movies \& TV ratings; see dataset documentation for citation details.]

\bibitem{libs} Python libraries: Pandas, NumPy, Matplotlib, SciPy.

\end{thebibliography}

\end{document}
